\begin{textsample}{POS Dim 3 – human – Score 7.00 – es/es\_ef\_anos\_iniciais\_linguagens\_lp\_2\_p3.txt}  \label{ex:f3_pos_004}
Feita essa reflexão, apresentamos esta \textbf{proposta} \textbf{curricular}, que se sistematiza considerando as competências e habilidades básicas \textbf{comuns} a serem consolidadas pelos estudantes na disciplina de Língua Portuguesa, ao longo dos diferentes momentos do Ensino Fundamental, desde o Ciclo da Alfabetização.
Também há, neste \textbf{documento}, sugestões para a operacionalização do trabalho docente no desenvolvimento das ações \textbf{educativas} com foco nas experiências escolares, \textbf{visto} que ao componente Língua Portuguesa \textbf{cabe} proporcionar aos estudantes experiências que contribuam para a ampliação do processo de aquisição das práticas da leitura e da escrita, de forma a possibilitar a participação significativa e crítica nas diversas práticas sociais permeadas/constituídas pela oralidade, pela escrita e por todas as formas de linguagens.

% matched lemmas: caber, comum, curricular, documento, educativo, proposta, ver
\end{textsample}
